% Generated draft: requires mathematical and source review.
\documentclass[11pt]{article}
\usepackage[margin=25mm]{geometry}
\usepackage{fontspec}
\setmainfont{Latin Modern Roman}
\usepackage{amsmath,amssymb,mathtools,unicode-math}
\setmathfont{Latin Modern Math}
\usepackage{enumitem,needspace,xurl,hyperref}
\setlength{\parindent}{0pt}
\setlength{\parskip}{5pt}
\setlength{\emergencystretch}{3em}
\newcommand{\N}{\mathbb{N}}
\newcommand{\Z}{\mathbb{Z}}
\newcommand{\Q}{\mathbb{Q}}
\newcommand{\R}{\mathbb{R}}
\newcommand{\C}{\mathbb{C}}
\newcommand{\F}{\mathbb{F}}
\newcommand{\A}{\mathbb{A}}
\newcommand{\PP}{\mathbb{P}}
\newcommand{\E}{\mathbb{E}}
\newcommand{\eps}{\varepsilon}
\newcommand{\dd}{\,\mathrm{d}}
\begin{document}
\section*{UnsolvedMath ID: 20000834}
\textbf{Source problem code:} catalogue entry
\subsection{Definitions and mathematical statement}
Use the full Hilbert scheme of subschemes of $\PP^3$ with polynomial $dt+1-g$, allowing embedded points; it cannot be restricted to the locally Cohen--Macaulay open locus. An embedded point is a zero-dimensional associated point lying on the one-dimensional support, distinct from an isolated point off that support.

\textbf{Problem 33.} Find the best pair $(d,g)$ for which an irreducible component has a general member with an embedded point, with $d$ as small and $g$ as close to zero as possible. The source expressly calls this optimization vague and supplies no ordering of the competing goals. Its known example is $(d,g)=(4,-15)$, cited to Chen--Nollet; that is an achieved pair, not a proved optimum.

\textbf{Reported detachment obstruction.} In the projective-space/field setting of the cited detachment theorem, the aid states that if a general member is obtained from a codimension-two locally complete-intersection curve by embedded points whose local quotient lengths are all at most three, those points detach in the same irreducible component. Consequently this situation cannot produce an embedded-general component. Here locally complete intersection means that the curve ideal is locally generated by a regular sequence of length two, and local multiplicity refers to length of the added zero-dimensional quotient.

Thus a candidate embedded-general component must use either a non-LCI underlying curve or an embedded local quotient of length at least four. This is a necessary branch division, not a sufficiency theorem for either branch.

To make the genus bookkeeping explicit, an embedded extension is described by
\[
 0\longrightarrow\mathcal I_Z\longrightarrow\mathcal I_C\longrightarrow Q\longrightarrow0,
 \qquad\ell=\operatorname{length}Q,
\]
so $\deg Z=\deg C=d$ and $p_a(Z)=p_a(C)-\ell$. In the LCI branch, $\ell\geq4$ and the usual degree-$d$ curve genus bound used by the report gives the necessary screen
\[
 g\leq\binom{d-1}{2}-4.
\]
This screen does not prove that an appropriate non-detaching extension or component exists.

\textbf{Reported simple-extension dimension threshold.} At a point $p$ where $\mathcal I_C$ has $r$ minimal generators, a simple quotient $\mathcal I_C\twoheadrightarrow k(p)$ is parametrized by $\PP^{r-1}$. If such support points move in a one-dimensional stratum of $C$, the extension locus has relative dimension $1+(r-1)=r$, whereas one detached point moves in three dimensions in $\PP^3$. Thus $r=3$ is the first dimensionally competitive case in this comparison. If the $r$-generator stratum is isolated, the extra one-dimensional term is absent; no unrestricted dimension formula is inferred.

The aid says the known thick-quartic-line example realizes the three-generator threshold. It supplies neither the explicit example ideal nor a complete component classification. The threshold is a search guide and the $(4,-15)$ example does not determine the source's optimum. The precise characteristic, detachment hypotheses and component-level argument require checking in Chen--Nollet.
\subsection{Short English statement}
The source seeks the smallest-degree, closest-to-zero-genus component whose general curve has an embedded point. The aid gives necessary obstructions: low-length embedded extensions of LCI curves detach, while a simple moving extension first competes dimensionally at three generators. The known $(4,-15)$ example is retained without claiming optimality.
\subsection{Sources}
\begin{itemize}
\item Local UnsolvedMath registry, permanent ID $20000834$, \texttt{catalogue entry}; canonical \texttt{aim-arithmetic-geometry-notes.json}, \texttt{notes[79]}.
\item AIM workshop ``Components of Hilbert schemes,'' PDF Problem 33 and nested remark: \url{https://aimath.org/WWN/hilbertschemes/hilbertschemes.pdf}.
\item D. Chen, S. Nollet, ``Detaching embedded points,'' \url{https://arxiv.org/abs/0911.2221}, cited in the original for $(d,g)=(4,-15)$.
\item Research attempt 1, \texttt{partial\_result}, run status \texttt{open}, internal rigor \texttt{proved}, low novelty confidence; current registry partially solved. Explicitly machine-generated, not peer reviewed, and to be independently verified before citation. Original literature field empty.
\end{itemize}
\end{document}
